\section{The Agora System}
\label{sec:method}

\subsection{World Representation}
Given a one-sentence premise $p$, the generation pipeline constructs
\begin{equation}
W=(E,G,R,K,A,M,P),
\end{equation}
where $E$ contains entities and their stable identities; $G$ is traversable
geometry; $R$ contains physical, economic, and institutional rules; $K$
specifies knowledge allocation and access; $A$ is the initial action
repertoire; $M$ associates visual media with entities; and $P$ records world
identity and revision. Mutable state $S_t$ contains locations, holdings,
knowledge acquired during play, tasks, relationships, and event status.
The schema is shared across worlds; its entities and rules depend on $p$.

This separation gives language-generated descriptions operational meaning.
An archive must refer to an existing room, an item must have a valid owner,
and a private fact must have a defined holder. Each visual asset is linked to
the same entity identity used by the runtime. Changes then propagate through
the state observed by subsequent actors rather than remaining only in dialogue.

\begin{figure*}[t]
\centering
\includegraphics[width=\textwidth]{figures/generation_overview.pdf}
\caption{\paneltitle{From premise to persistent consequences.} Specialized
generation stages populate a shared world representation. Validation connects
semantic, visual, and spatial requirements; local repair revisits affected
dependencies. Humans and agents act through the same world, while logged
decisions and state changes support separate construction and interaction
evaluations.}
\label{fig:generation-overview}
\end{figure*}

\subsection{Compositional Generation and Local Refinement}
A planning stage identifies institutions, conflicts, resources, and activity
loops. Specialized stages expand materials, rooms, items, roles, main
characters, knowledge, visual direction, wardrobe, property, and placement.
The shared representation constrains references between their outputs.
Visual generation produces room floors, character atlases, and objects under
a common style specification; assembly connects these assets to the playable
geometry and entity identities.

Validation targets four properties: semantic completeness, cross-modal
identity, navigable and readable space, and persistent interactive state.
Structural validators check references, inventories, asset provenance, and
image contracts. Visual inspection and interactive checks supply different
evidence about readability and use. When a check identifies an invalid
character--item relation, that dependency can be revisited while unaffected
stages remain fixed. This instantiates repair as a constrained part of
generation, consistent with the broader PCG literature \cite{pcgml}.

\paragraph{Running example.}
Tidal Embassy of Lost Languages turns linguistic scarcity into institutions
and actions. Its generated configuration contains eight rooms, 25 named main
characters, and 20 item types. Diplomats, smugglers, and memory divers receive
different possessions and private knowledge. The setting exposes activities
such as translating an archaic bribe, forging a linguistic seal, and rescuing
a flooded treaty. These relations create reasons to move, exchange, withhold
information, and seek cooperation. Supplementary Section~\ref{supp-app:tidal} follows the full
generation trace, including failed room images and localized visual repair.

\subsection{Bounded Observation and Situated Intention}
At round $t$, an activation policy selects agents. Agent $i$ receives
\begin{equation}
o_i^t=(\ell_i^t,N_i^t,I_i^t,K_i^t,H_i^t,Q_i^t),
\end{equation}
comprising location, nearby actors, inventory, accessible knowledge, episode
and task memory, and selected directed relationships. Role and goals provide
the decision context. The observation boundary is intended to preserve
information asymmetry, an important distinction from omniscient social
simulation \cite{misleading}.

The agent chooses a partner and an intention, using either a generated action
or an open proposal. Visible actions explain what an institution makes
possible; proposals express intentions absent from that initial vocabulary.
Open proposals in the audited implementation support four effect primitives:
speech, transfer, item creation, and temporary status. Movement is a separate
action route. These primitives can describe new contextual intentions, but
they are not arbitrary executable programs or a general physics engine. Humans can approach an inhabitant, send messages, use or trade
items, select a visible action, or formulate a proposal through the same world
interface. This extends the shared grounded-play tradition of LIGHT
\cite{light} to premise-generated settings.

\subsection{Coordinator-Governed State Changes}
We write an open proposal as $q=(i,j,u,F,C)$: actor $i$, optional target $j$,
language intention $u$, requested effects $F$, and declared consent evidence
$C$. The model supplies a structured proposal; the audited coordinator core uses
deterministic schema and state checks to return a decision and explanation,
\begin{equation}
\Gamma(W,S_t,q)=(d,r),\quad
d\in\{\text{approve},\text{revise},\text{reject}\}.
\end{equation}
Following the action-grounding role of Concordia's Game Master
\cite{vezhnevets2023}, interpretation is separated from committing effects.
An approved effect is applied by the state-transition machinery. Revision or
rejection leaves the requested material effects unapplied, while the attempt
can remain in the interaction history.

\begin{table}[t]
\centering\small
\setlength{\tabcolsep}{4pt}
\begin{tabularx}{\columnwidth}{@{}L{0.19\columnwidth}Y@{}}
\toprule
Effect & Required grounding \\
\midrule
Speech & Addressable participant and a permitted channel \\
Transfer & Existing item, ownership, quantity, and co-presence where required \\
Trade & Transfer conditions and the partner's acceptance \\
Creation & Declared available inputs and a resulting object identity \\
Status & Permitted target, bounded magnitude, and duration \\
Movement & Reachable destination and valid occupancy \\
\bottomrule
\end{tabularx}
\caption{Examples of action conditions in the shared world. Linguistic
plausibility alone does not establish that a requested effect is executable.}
\label{tab:conditions}
\end{table}

The resulting loop combines generative interpretation with state constraints
(Table~\ref{tab:conditions}). Ownership, quantity, location, and duration are
represented explicitly. Social interpretation, including the meaning of an
exchange or an accusation, remains model-mediated upstream. Declared partner
acceptance is an input to validation; it is not itself proof of independently
obtained human consent. Natural-language institutions are not automatically
formalized as executable invariants by these checks. Consequently, the
coordinator is a practical adjudication mechanism, not a proof of correct
interpretation or universally valid social norms. Our evaluation reports
observed outcomes rather than claiming exhaustive rule compliance.

\begin{figure*}[t]
\centering
\includegraphics[width=\textwidth]{figures/situated_action_cycle.pdf}
\caption{\paneltitle{One interaction cycle.} Bounded observations inform a
visible action or an open proposal. The coordinator evaluates the proposed
effects; approved changes update shared state. Both outcomes and failed
attempts provide history for later decisions.}
\label{fig:agent-loop}
\end{figure*}

\section{Evaluation Design}
\label{sec:evaluation}
We evaluate three links in the proposed mechanism: generated entities must
satisfy shared representation contracts; inhabitants must be able to use both
discoverable actions and contextual proposals; approved effects must persist
while invalid effects are blocked. We separate generated worlds, interaction
events, action results, and participant judgments as different units of evidence. A world can pass structural checks without producing
engaging play. An event can contain several effects or attempts, so the
number of logged events need not equal the number of adjudicated results.
Longer trajectories describe the simulated population, while human judgments
require a separate instrument and participants.

\begin{table*}[t]
\centering\small
\setlength{\tabcolsep}{5pt}
\begin{tabularx}{\textwidth}{@{}L{0.19\textwidth}L{0.23\textwidth}YY@{}}
\toprule
Study & Unit and scope & Evidence reported & Interpretation \\
\midrule
Internal generation comparison & 10 paired premises; 20 trials per method & Completion, structural quality proxies, cost & Descriptive reliability--cost tradeoff \\
Validation and repair & 58 injected faults with clean checks; 5 repair cases & Detection, attribution, local repair cost & Specified contracts and seeded defects \\
Multiworld interaction & 5 worlds; 2 rounds each; 53 events & Action pathways, dyads, outcomes, state changes & Initial behavior under common runtime settings \\
Cross-model audit & 24 generations; 21 scripted 24-round runs & Completion, executed effects, rejected actions & Retrospective coverage and execution stress test \\
Historical pilot & 3 prior runs; 1,052 events & Repeated ties and relational trajectories & Exploratory, confounded across worlds \\
Human model comparison & 10 eligible returns; 40 image comparisons; 7 configurations & Exploration, map readability, character fit & Fixed pipeline and premise; sparse artifact preferences \\
Internal procedure control & 20 structure comparisons; 10 requests & Same-system generation variants; full counts in supplement & Unequal budgets; no external-system comparison \\
Supplementary planner diagnostics & 18 trajectories; 2 fixed worlds; 3 configurations & Input masking, available effects, and state replay & One centralized planner; implementation-specific counters \\
\bottomrule
\end{tabularx}
\caption{Evaluation scope and evidence status. The studies distinguish construction, artifact preference, and
executed consequences; evidence from different protocols is not pooled.}
\label{tab:evidence}
\end{table*}

The evaluation addresses three mechanism questions: \textbf{RQ1}, how do
shared contracts and local repair connect generated components into an
executable world; \textbf{RQ2}, how are generated action opportunities and
contextual proposals exercised in that world; and \textbf{RQ3}, how do
coordinator decisions become persistent state changes, and what happens to
invalid effects? Construction checks address RQ1, the five-world interaction
study addresses RQ2, and execution audits with replayable traces address RQ3.
The completed visual study characterizes the resulting artifacts. Internal
generation controls and supplementary planner ablations report operating
conditions and implementation boundaries. Table~\ref{tab:evidence} states the
unit and interpretation of each study.
